% p.79 — الباب الرابع والثلاثون في قوانين التعليل وغيره (تابع: القوانين 21–31)
% Layout: twocol. Right column: end of rule 21 («لان الحذف …»), rules 22–25 and the start of 26
% (ends «اِحْذِفْ تاء المخاطبة»); left column: rest of 26 («و بعده سكون …»), 27–31.
\mspage{79}
\begin{twocol}
\begin{rightcol}
لان الحذف هو علامة النصب والجزم
\par
(22) كذالك تقول فى الناقص لَمْ يَغْزُوَا لَمْ يَغْزُوْا لَمْ يَرْمِيَا لَمْ يَرْمُوْا لَمْ يَخْشَيَا لَمْ يَخْشَوْا بحذف النون ولا تحذف الواو والياء والالف كما حذفت فى المفرد
\par
(23) كذالك تقول فى الاجوف لَمْ يَقُوْلَا لَمْ يَقُوْلُوْا لَمْ يَبِيْعَا لَمْ يَبِيْعُوْا لَمْ يَخَافَا لَمْ يَخَافُوْا بحذف النون لكن لا تحذف الواو والياء والالف كما حذفت فى المفرد لانه لَمْ يَلْتَقِ السَّاكِنَانِ
\par
(24) لا تحذف نون الجمع المؤنث قط تقول يَضْرِبْنَ تَضْرِبْنَ لم يَضْرِبْنَ لَنْ يَضْرِبْنَ لانه مَبْنِيٌّ اَبَدًا
\par
(25) الحروف كلها مبنية والافعال كذالك الا الفعل المضارع ولا يدخل الجزم فى الاسماء والكسرةُ فى الافعال
\par
(26) اعلم ان الامر الحاضر يؤخذ من المضارع المخاطب مثلا تقول فى تَعِدُ عِدْ احذف تاء المخاطبة و اسكن الآخر وتقول فى تَضْرِبُ و تَفْتَحُ اِضْرِبْ اِفْتَحْ اِحْذِفْ تاء المخاطبة
\end{rightcol}
\begin{leftcol}
و بعده سكون لا يمكن ان تَفْتَتِحَ به فاجلب همزة الوصل مكسورا و اذا كان العين مضموما فمضموما نحو تَنْصُرُ اُنْصُرْ
\par
(27) واذا كان المضارع ناقصا فتحذف الاخر علامة للسكون كما مر مثل تقول فى تَغْزُوْ اُغْزُ وفى تَرْمِىْ اِرْمِ فى تَخْشٰى اِخْشَ وفى \com{التثنة}{كذا في الأصل، ولعل الصواب: التثنية} والجمع تقول اُغْزُوَا اِرْمِيَا اِخْشَيَا بحذف النون علامة للسكون بلا حذف الواو والالف والياء
\par
(28) حكم همزة الوصل ان تثبت فى الابتداءِ وتسقط فى الدرج نحو تقول اِضْرِبْ ابتداء مع الهمزة واذا وصلت مع كلمة فتسقط مثلا فَاضْرِبْ
\par
(29) يختم الكلمة بالسكون فيسمّى وقفا تقول زَيْدْ عَمْرْو واذا كان فى اخر الكلمة فتحة منونا فبالالف فتقول فى زيدًا زَيْدَا \com{واذ كان}{كذا في الأصل، ولعل الصواب: واذا كان} تاءً فبالهاء كما تقول فى غاشيةٍ غاشِيَهْ
\par
(30) واذا كان الكلمة حرفا واحدا فتزيد هاء السكتة عِنْدَ الوقف فتقول فى قِ قِهْ
\par
(31) ان الاصل فى تُكْرِمُ تُأَكْرِمُ فتقول فى الامر اَكْرِمْ بعد حذف تاء الخطاب مع اسكان الآخر
\end{leftcol}
\end{twocol}
